\documentclass[10pt,a4paper]{amsart}
\usepackage[margin=19mm]{geometry}
\usepackage{amsmath,amssymb,mathtools}
\usepackage[scr=rsfs,cal=euler]{mathalfa}
\usepackage{xfrac}
\usepackage[unicode,hidelinks,bookmarks=false]{hyperref}
\newcommand{\ie}{i.e.\ }
\newcommand{\cf}{cf.\ }
\newcommand{\resp}{respectively\ }
\newcommand{\Subsection}{\subsection}
\providecommand{\nobreakdash}{\nobreak\hbox{-}}
\setlength{\emergencystretch}{2em}
\multlinegap0pt
\usepackage{bm}
\usepackage{bbm}
\usepackage{dsfont}
\usepackage{xspace}
\usepackage{cancel}
\usepackage{mathtools}
\usepackage{amsthm}
\makeatletter
\@ifundefined{Lemma}{\newtheorem{Lemma}{Lemma}}{}
\@ifundefined{Proposition}{\newtheorem{Proposition}[Lemma]{Proposition}}{}
\makeatother
\newcommand{\R}{\mathbb{R}}
\newcommand{\bC}{\mathsf{C}}
\newcommand{\bG}{\mathsf{G}}
\newcommand{\cP}{\mathscr{P}}
\newcommand{\comp}{\textup{cmp}}
\newcommand{\fg}{\mathfrak{g}}
\title[On the unramified Eisenstein spectrum]{On the unramified Eisenstein spectrum}
\author{David Kazhdan, Andrei Okounkov}
\date{Source: arXiv:2203.03486v3 (CC BY 4.0)}
\begin{document}
\maketitle

\noindent\emph{Source and licence.} On the unramified Eisenstein spectrum. Version arXiv:2203.03486v3 (CC BY 4.0), distributed under the Creative Commons Attribution 4.0 International licence, as recorded in the arXiv licence metadata for this exact version and shown on the abstract page https://arxiv.org/abs/2203.03486v3. Full rights notice: licenses/arxiv-2203.03486-CC-BY-4.0.txt.\par\smallskip

\noindent\textit{Editorial scope.} Counted unit: the Proposition whose source label is \texttt{p\_even} in the article (source0.tex lines 5609--5616, source label \texttt{p\_even}), delivered together with the article's own complete original proof of it (source0.tex lines 5618--5664, one \texttt{proof} environment of 47 lines). No mathematics is written, repaired or re-derived: statement and proof are the article's own text line for line. Required context retained uncounted: l\_selfdual (lines 5432--5439); eq:30 (lines 5489--5494); eq:27 (lines 5578--5583); eq:27bis (lines 5578--5583). Every definition, display and label of those lines is retained unchanged. Omitted from this package and not redistributed: 1--5431, 5440--5488, 5495--5577, 5584--5608, 5617--5617, 5665--8545. Nothing inside a retained excerpt is omitted or altered: every retained body is delivered exactly as the article's lines are written, so no omission receipt of any kind accompanies this package. Citations: the retained excerpts carry no citation command at all, so no bibliography entry of the article is transcribed and no reference is redistributed. Reference adaptation: none was needed and none was made; every reference inside the delivered unit resolves inside it, so no unresolved reference or citation is produced. Printed slips kept literal: no printed word, symbol, hypothesis or number of the counted statement or of its proof was altered. Layout and class: the article's own class is \texttt{extarticle} with packages including amssymb, amsfonts, amsthm, amsmath, bbold, mathrsfs, pifont, mathabx, xy, array, transparent, stmaryrd, tabularx, geometry; the wrapper is the lane's pinned \texttt{amsart} editorial wrapper at 10pt on A4, which restates the theorem-like environments the retained text needs and adds the article's own macro definitions that the retained text needs (\texttt{\textbackslash R}, \texttt{\textbackslash bC}, \texttt{\textbackslash bG}, \texttt{\textbackslash cP}, \texttt{\textbackslash comp}, \texttt{\textbackslash fg}), with extra wrapper packages (bm, bbm, dsfont, xspace, cancel, mathtools). The delivered numbering is the unit's own. Assets: no retained span contains a figure environment, a graphics-inclusion command, a PDF-page inclusion, a table or a TikZ picture; no image file is copied into this package, the package declares no asset dependency and no image file is redistributed with it. Licence: version arXiv:2203.03486v3 of this article is distributed under the Creative Commons Attribution 4.0 International licence, as recorded in the arXiv licence metadata for this exact version and shown on the abstract page https://arxiv.org/abs/2203.03486v3. A full rights notice is delivered at licenses/arxiv-2203.03486-CC-BY-4.0.txt. Characters: every retained excerpt is pure ASCII, so the delivered engine, XeLaTeX with its default Latin Modern text font, reports no missing character.
  \begin{Lemma}\label{l_selfdual}
    The $\bC_\phi$-modules $\fg_i$ and $\fg_i^e$ are selfdual.
    The $\bC_\phi$-module $\fg_1\cong \fg_{-1}$ is symplectic.
    We have
    \[
      \det \bigoplus_{i \ge 0} \fg_i^e = 1 \,. 
    \]
\end{Lemma}

    \begin{equation}
      \label{eq:30}
      Z_e \big|_{q^{h/2} c} = q^{-\frac{\dim \fg + \dim \fg^e}{2}} \frac{(1-q)^{2r}}{Z^1(\fg)}
      \frac{\prod_{i>0} \det_{\fg_i^e} (1-q^{-i/2} c )\phantom{\,\,\,\,\,}}
      {\prod_{i\ge 0} \det_{\fg_i^e}(1-q^{-i/2-1} c)} \,, \quad c \in \bC_{\phi} \,. 
    \end{equation}

 \begin{align}
   q &>1 \,, \notag \\
   Z(x)&> 0\,, \quad x>q \,,  \label{eq:27} \\
   Z^1(x)&\ne 0\,,\quad  x \in \R \setminus \{1 < |x| < q\} \label{eq:27bis}
         \,. 
 \end{align}

 \begin{Proposition}\label{p_even} 
   With the assumptions \eqref{eq:27} we have
   \[
     \cP_{-}^Z f^* \Big|_{q^{h/2} \bC_{\phi,\comp}} =
   \overline{\cP_{+}^Z f} \Big|_{q^{h/2} \bC_{\phi,\comp}}
   \]
 and $Z_e \ge 0$ on $q^{h/2} \bC_{\phi,\comp}$. 
\end{Proposition}

\begin{proof}
Since the function $Z(x)$ is single-valued and real on a part of the
real axis, it follows that it is real, that is, $\overline{Z(x)} =
Z(\overline{x})$.

Consider an element of the form
\[
  x = q^{h/2} c \,, \quad c \in \bC_{\phi,\comp} \,, 
\]
and note that
\[
  \overline{x} =  q^{h/2} \overline{c} \sim q^{h/2} c^{-1} \sim 
q^{-h/2} c^{-1} = x^{-1} \,, 
\]
where similarity means being conjugate in $\bG$. Indeed,
$\overline{c}\sim c^{-1}$ inside a maximal compact subgroup of
$\bG^h$, while $h$ is conjugate to
$-h$ by the image of $\phi: SL(2)\to \bG^c$. 
Since $\cP_{+}^Z f$ is conjugation-invariant, we conclude 
\[
  \overline{\cP_{+}^Z f} \Big|_{x=q^{h/2}c }= \cP_{+}^Z \overline{f} \Big|_{x^{-1}} =
\cP_{-}^Z f^* \Big|_{x} \,.
\]

We now consider $Z_e$. Since eigenvalues of $c$ lie on the unit circle
in complex conjugate pairs, we observe that all terms in \eqref{eq:30}
are positive, with the possible exception of $Z^1(\fg)$.

The
representations $\fg_i$ being self-dual, its $c$-weights appear either
in complex conjugate pairs or are equal to $\pm 1$. We are thus left
with the contributions of $Z^1(\pm q^{i/2})$.

By our assumption, $Z^1(q^{i/2})>0$, except possibly $Z^1(q^{1/2})$.
However, the
representation $\fg_1$ being symplectic by Lemma \ref{l_selfdual}, the
$q^{1/2}$-weights in $\fg_1$ appears with even multiplicity, hence
its contribution is positive. 

Now consider the contribution of $Z^1(- q^{i/2})$. By \eqref{eq:27bis} all of these
numbers are of the same sign, except possibly $Z^1(- q^{1/2})$, which
appears with even multiplicity in the symplectic representation
$\fg_1$. The nontrivial weights in the adjoint
representation appear in pairs given by the opposite roots. In
particular, the total number of negative real weights in $\fg$ is even for
any element of $\bG$. This completes the proof. 
\end{proof}

\end{document}
